% GENERATED FILE. Edit canonical lessons, then run npm run generate:books.
% canonical-source-hash: fnv1a64:d2983f6a3db49913
% canonical-lessons: GE-C284-verdacht, GE-C284-verdaechtig, GE-C284-ermitteln, GE-C284-ermittlung, GE-C284-festnehmen

\chapter{Der Verdacht}
\label{ch:ge-der-verdacht}
\hlchaptermodality{284}

\begin{chapteropening}
\textbf{By the end of this chapter:} \emph{I can talk about an investigation: suspicion, suspicious, to investigate, an investigation and to arrest.}

Complete the last lesson of chapter 284: I can talk about an investigation: suspicion, suspicious, to investigate, an investigation and to arrest.
\end{chapteropening}

\section[der Verdacht]{der Verdacht — the suspicion}

\label{lesson:GE-C284-verdacht}



\begin{quote}
Before the new one: say the German for to observe, then the German for to recount.
\end{quote}

\subsection*{You'll want to know: der Verdacht}
\textbf{der Verdacht} — \textquotedblleft{}the suspicion\textquotedblright{}.

\begin{quote}
Der Verdacht fiel auf den Nachbarn.
\end{quote}

\begin{quote}
\emph{Suspicion fell on the neighbour.}
\end{quote}

No plural. \emph{Unter Verdacht stehen} — to be under suspicion.

\begin{cousinweb}
\emph{Verdacht} is usually traced to the old verb \emph{verdenken} \textquotedblleft{}to suspect\textquotedblright{}, built on \emph{denken} \textquotedblleft{}to think\textquotedblright{}; EtymDB lists no parent.
\end{cousinweb}

\begin{grammarlens}[title={What you've built}]
One more word for telling, explaining and describing at B1.
\end{grammarlens}

\subsection*{Your turn}
\begin{itemize}
\raggedright
  \item \emph{Say it:} \emph{der Verdacht}
  \item \emph{Say it:} \emph{der Verdacht}, once more
  \item \emph{Say it:} a sentence of your own with \emph{der Verdacht}
  \item \emph{Recall:} say the German for to observe, then the German for to recount, then say \emph{der Verdacht} again
\end{itemize}

\begin{tcolorbox}[breakable,colback=teal!4,colframe=teal!35!black,title={Before you move on}]
What does der Verdacht mean? (\textquotedblleft{}The suspicion\textquotedblright{}.) Say it once more.
\end{tcolorbox}
\section[verdächtig]{verdächtig — suspicious}

\label{lesson:GE-C284-verdaechtig}



\begin{quote}
Before the new one: say the German for to recount, then the German for the suspicion.
\end{quote}

\subsection*{You'll want to know: verdächtig}
\textbf{verdächtig} — \textquotedblleft{}suspicious\textquotedblright{}.

\begin{quote}
Ein verdächtiges Auto stand vor dem Haus.
\end{quote}

\begin{quote}
\emph{A suspicious car stood in front of the house.}
\end{quote}

\emph{Der Verdächtige} — the suspect.

\begin{cousinweb}
\emph{verdächtig} is built on \emph{Verdacht} \textquotedblleft{}suspicion\textquotedblright{}, which is usually traced to \emph{verdenken}, from \emph{denken} \textquotedblleft{}to think\textquotedblright{}; EtymDB lists no parent.
\end{cousinweb}

\begin{grammarlens}[title={What you've built}]
One more word for telling, explaining and describing at B1.
\end{grammarlens}

\subsection*{Your turn}
\begin{itemize}
\raggedright
  \item \emph{Say it:} \emph{verdächtig}
  \item \emph{Say it:} \emph{verdächtig}, once more
  \item \emph{Say it:} a sentence of your own with \emph{verdächtig}
  \item \emph{Recall:} say the German for to recount, then the German for the suspicion, then say \emph{verdächtig} again
\end{itemize}

\begin{tcolorbox}[breakable,colback=teal!4,colframe=teal!35!black,title={Before you move on}]
What does verdächtig mean? (\textquotedblleft{}Suspicious\textquotedblright{}.) Say it once more.
\end{tcolorbox}
\section[ermitteln]{ermitteln — to investigate; to find out}

\label{lesson:GE-C284-ermitteln}



\begin{quote}
Before the new one: say the German for the suspicion, then the German for suspicious.
\end{quote}

\subsection*{You'll want to know: ermitteln}
\textbf{ermitteln} — \textquotedblleft{}to investigate; to find out\textquotedblright{}.

\begin{quote}
Die Polizei ermittelte wochenlang.
\end{quote}

\begin{quote}
\emph{The police investigated for weeks.}
\end{quote}

\emph{Gegen jemanden ermitteln} — to investigate someone.

\begin{cousinweb}
\emph{ermitteln} is built on \emph{Mittel} \textquotedblleft{}middle, means\textquotedblright{}, an inherited Germanic word (EtymDB), with the prefix \emph{er-}; EtymDB lists no parent for the verb itself.
\end{cousinweb}

\begin{grammarlens}[title={What you've built}]
One more word for telling, explaining and describing at B1.
\end{grammarlens}

\subsection*{Your turn}
\begin{itemize}
\raggedright
  \item \emph{Say it:} \emph{ermitteln}
  \item \emph{Say it:} \emph{ermitteln}, once more
  \item \emph{Say it:} a sentence of your own with \emph{ermitteln}
  \item \emph{Recall:} say the German for the suspicion, then the German for suspicious, then say \emph{ermitteln} again
\end{itemize}

\begin{tcolorbox}[breakable,colback=teal!4,colframe=teal!35!black,title={Before you move on}]
What does ermitteln mean? (\textquotedblleft{}To investigate; to find out\textquotedblright{}.) Say it once more.
\end{tcolorbox}
\section[die Ermittlung]{die Ermittlung — the investigation}

\label{lesson:GE-C284-ermittlung}



\begin{quote}
Before the new one: say the German for suspicious, then the German for to investigate.
\end{quote}

\subsection*{You'll want to know: die Ermittlung}
\textbf{die Ermittlung} — \textquotedblleft{}the investigation\textquotedblright{}.

\begin{quote}
Die Ermittlungen dauerten zwei Monate.
\end{quote}

\begin{quote}
\emph{The investigations took two months.}
\end{quote}

Mostly plural: \emph{die Ermittlungen}.

\begin{cousinweb}
\emph{Ermittlung} is built on \emph{ermitteln} \textquotedblleft{}to investigate\textquotedblright{}, from \emph{Mittel} \textquotedblleft{}middle, means\textquotedblright{} (EtymDB), with the ending \emph{-ung}.
\end{cousinweb}

\begin{grammarlens}[title={What you've built}]
One more word for telling, explaining and describing at B1.
\end{grammarlens}

\subsection*{Your turn}
\begin{itemize}
\raggedright
  \item \emph{Say it:} \emph{die Ermittlung}
  \item \emph{Say it:} \emph{die Ermittlung}, once more
  \item \emph{Say it:} a sentence of your own with \emph{die Ermittlung}
  \item \emph{Recall:} say the German for suspicious, then the German for to investigate, then say \emph{die Ermittlung} again
\end{itemize}

\begin{tcolorbox}[breakable,colback=teal!4,colframe=teal!35!black,title={Before you move on}]
What does die Ermittlung mean? (\textquotedblleft{}The investigation\textquotedblright{}.) Say it once more.
\end{tcolorbox}
\section[festnehmen]{festnehmen — to arrest, to detain}

\label{lesson:GE-C284-festnehmen}



\begin{quote}
Before the new one: say the German for to investigate, then the German for the investigation.
\end{quote}

\subsection*{You'll want to know: festnehmen}
\textbf{festnehmen} — \textquotedblleft{}to arrest, to detain\textquotedblright{}.

\begin{quote}
Die Polizei nahm den Mann am Bahnhof fest.
\end{quote}

\begin{quote}
\emph{The police arrested the man at the station.}
\end{quote}

Separable: \emph{sie nehmen ihn fest}, \emph{sie haben ihn festgenommen}. \emph{Verhaften} means the same.

\begin{cousinweb}
\emph{festnehmen} joins \emph{fest} \textquotedblleft{}firm\textquotedblright{}, Old High German \emph{festi} (EtymDB), and \emph{nehmen} \textquotedblleft{}to take\textquotedblright{}, Proto-Germanic \emph{nemaną} (EtymDB): to take and hold.
\end{cousinweb}

\begin{grammarlens}[title={What you've built}]
One more word for telling, explaining and describing at B1.
\end{grammarlens}

\subsection*{Your turn}
\begin{itemize}
\raggedright
  \item \emph{Say it:} \emph{festnehmen}
  \item \emph{Say it:} \emph{festnehmen}, once more
  \item \emph{Say it:} a sentence of your own with \emph{festnehmen}
  \item \emph{Recall:} say the German for to investigate, then the German for the investigation, then say \emph{festnehmen} again
\end{itemize}

\begin{tcolorbox}[breakable,colback=teal!4,colframe=teal!35!black,title={Before you move on}]
What does festnehmen mean? (\textquotedblleft{}To arrest, to detain\textquotedblright{}.) Say it once more.
\end{tcolorbox}
